\documentclass[12pt]{article}
\usepackage[utf8]{inputenc}
\usepackage[T1]{fontenc}
\usepackage{hyperref}
\title{Foma Language Resources for K-12 Curriculum Development}
\author{Compiled by Claude Sonnet 5.5 (claude-sonnet-5-5) for Andr\'as Kornai}
\date{}
\begin{document}
\maketitle

\section{Introduction}
Foma (Lifoma; ISO 639-3 \texttt{fom}) is spoken along the north bank of the Congo River upstream from Basoko, Tshopo Province, Democratic Republic of the Congo. Searches found very little dedicated material. Some sources treat Foma as a dialect of Kele-Foma (Lokele), while others list it as a separate language. This report is therefore short, and no school or revitalization materials were found.

\section{Primary Educational Resources}
No primers, textbooks or school materials in Foma were found in search results. Joshua Project reports that Lingala and French are used for church life, trade and education in the region.

\section{Linguistic and Reference Materials}
\begin{itemize}
\item Wikipedia, Kele-Foma language (general overview, classification, and the relation of Foma to Lokele): \url{https://en.wikipedia.org/wiki/Kele-Foma_language}
\end{itemize}
No dictionary or grammar of Foma was found.

\section{Religious and Community Resources}
\begin{itemize}
\item Joshua Project, Foma in Congo, Democratic Republic of (people-group profile; states that no Scripture resources in Foma are documented): \url{https://joshuaproject.net/people_groups/11221/CG}
\end{itemize}

\section{Digital and Audio Resources}
\begin{itemize}
\item Wikitongues, Foma (language page for \texttt{fom}; it may link to speaker recordings, but this was not verified): \url{https://wikitongues.org/languages/fom/}
\end{itemize}

\section{Curriculum Development Context}
Little is documented. Any curriculum would need new documentation work with the Foma community, and the relation to Lokele (Kele-Foma) should be clarified first. The Congolese school system uses French and Lingala, so a Foma programme would need local partners. Nothing here suggests existing vernacular schooling.

\end{document}